\section{Figures, Tables and Displayed Mathematics}
\label{sec:elements}

The three element types below cover most reports. They are shown once, correctly, so that a new
author can copy the pattern rather than reconstruct it.

\subsection{Figures}

Figures are centred, float by default, and carry a caption below. Vector output is expected:
a PDF produced by a plotting tool, or a diagram drawn in Ti\emph{k}Z as in
Figure~\ref{fig:pipeline}. Raster images are acceptable only for screenshots, and then at a
resolution that survives print.

\begin{figure}[htbp]
  \centering
  \begin{tikzpicture}[
      node distance=9mm,
      stage/.style={draw, rounded corners=2pt, semithick, minimum height=9mm,
                    minimum width=27mm, align=center, font=\small},
      arrow/.style={-{Latex[length=2mm]}, semithick}
    ]
    \node[stage] (draft)  {Draft};
    \node[stage, right=of draft]  (review) {Boogon\\review};
    \node[stage, right=of review] (typeset) {Typeset};
    \node[stage, right=of typeset] (release) {Release};
    \draw[arrow] (draft) -- (review);
    \draw[arrow] (review) -- (typeset);
    \draw[arrow] (typeset) -- (release);
  \end{tikzpicture}
  \caption{The release path for a Cooper technical report. Editorial review sits between
  drafting and typesetting, so that a report is never typeset twice for the same objection.}
  \label{fig:pipeline}
\end{figure}

\subsection{Tables}

Tables use \texttt{booktabs} rules and no vertical lines; the caption goes above. Numbers are
aligned on the decimal, and units are stated in the column head, never repeated in every cell.
Table~\ref{tab:conformance} is the smallest complete example.

\begin{table}[htbp]
  \centering
  \caption{Conformance of a report to the template, by section, for a recent review cycle.}
  \label{tab:conformance}
  \begin{tabular}{@{}lrr@{}}
    \toprule
    Section & Checks run & Conformant (\%) \\
    \midrule
    Front matter        & 14 & 100.0 \\
    Body sections       & 63 & 96.8 \\
    Figures and tables  & 22 & 90.9 \\
    References          &  9 & 100.0 \\
    \midrule
    Total               & 108 & 96.3 \\
    \bottomrule
  \end{tabular}
\end{table}

\subsection{Displayed mathematics}

Equations are numbered, and every symbol is defined in the surrounding prose. A conformance
rate, for instance, can be written as
\begin{equation}
  \label{eq:conformance}
  \mathcal{C} = \frac{1}{N}\sum_{i=1}^{N} \mathbf{1}\!\left[\, d_i \le \tau \,\right],
\end{equation}
where $d_i$ is the deviation of check $i$ from its expected value, $\tau$ is the tolerance
agreed during review, and $\mathbf{1}[\cdot]$ is the indicator function. Equation
\eqref{eq:conformance} is quoted in the review summary; it is also the quantity that
Table~\ref{tab:conformance} reports in percent.

Inline mathematics follows the same font. A tolerance of $\tau = 0.05$ is typical for
typographic checks. The treatment of floats and cross-references used here follows the
standard account in \citet{lamport1994}.
